% !TEX program = xelatex
% Lernplatt - by Can Siebert
% Kurs 22 (Bo 143) - Tag 12 - Aufgabenheft
% Revenue Streams optimieren, Multi-Channel skalieren, Wachstum steuern
%
% Self-contained xelatex source.

\documentclass[11pt,a4paper,oneside]{article}

\usepackage{polyglossia}
\setdefaultlanguage{german}

\usepackage[a4paper,margin=2.5cm]{geometry}

\usepackage{fontspec}
\defaultfontfeatures{Ligatures=TeX}

\IfFontExistsTF{Charter}{%
  \setmainfont{Charter}%
}{%
  \IfFontExistsTF{Noto Serif}{%
    \setmainfont{Noto Serif}%
  }{%
    \setmainfont{Liberation Serif}%
  }%
}

\IfFontExistsTF{Source Sans 3}{%
  \setsansfont{Source Sans 3}%
}{%
  \IfFontExistsTF{Source Sans Pro}{%
    \setsansfont{Source Sans Pro}%
  }{%
    \setsansfont{Liberation Sans}%
  }%
}

\newcommand{\caveatfontfallback}{\itshape}
\IfFontExistsTF{Caveat}{%
  \newfontfamily\caveatfont{Caveat}[Scale=1.15]%
}{%
  \newcommand\caveatfont{\caveatfontfallback}%
}

\setmonofont{Liberation Mono}

\usepackage{microtype}
\usepackage{eurosym}
\linespread{1.15}

\usepackage{xcolor}
\definecolor{lpInk}{RGB}{20,28,38}
\definecolor{lpAccent}{RGB}{22,90,140}
\definecolor{lpAccentLight}{RGB}{210,228,240}
\definecolor{lpAmber}{RGB}{222,138,30}
\definecolor{lpAmberLight}{RGB}{255,243,217}
\definecolor{lpAmberBorder}{RGB}{196,118,18}
\definecolor{lpGray}{RGB}{96,104,114}
\definecolor{lpGrayLight}{RGB}{238,240,243}
\definecolor{lpGreen}{RGB}{34,120,72}
\definecolor{lpGreenLight}{RGB}{222,238,228}
\definecolor{lpGreenBorder}{RGB}{24,96,58}

\definecolor{lpL1}{RGB}{34,120,72}
\definecolor{lpL2}{RGB}{22,90,140}
\definecolor{lpL3}{RGB}{170,42,52}

\usepackage[most]{tcolorbox}
\tcbuselibrary{breakable,skins}

\newtcolorbox{infobox}[1][Hinweis]{%
  enhanced, breakable,
  colback=lpAccentLight,
  colframe=lpAccent,
  boxrule=0.6pt,
  arc=2pt,
  left=10pt,right=10pt,top=8pt,bottom=8pt,
  fonttitle=\sffamily\bfseries\color{white},
  coltitle=white,
  title={#1},
  attach boxed title to top left={xshift=8pt,yshift=-8pt},
  boxed title style={size=small, colback=lpAccent, colframe=lpAccent, boxrule=0.4pt, arc=1pt},
  top=14pt
}

\newtcolorbox{antwortbox}[1][Ihre Antwort]{%
  enhanced, breakable,
  colback=white,
  colframe=lpGray,
  boxrule=0.4pt,
  arc=2pt,
  left=10pt,right=10pt,top=8pt,bottom=8pt,
  fonttitle=\sffamily\bfseries\color{white},
  coltitle=white,
  title={#1},
  attach boxed title to top left={xshift=8pt,yshift=-8pt},
  boxed title style={size=small, colback=lpGray, colframe=lpGray, boxrule=0.4pt, arc=1pt},
  top=14pt
}

\usepackage{enumitem}
\setlist{itemsep=2pt, topsep=4pt, parsep=0pt}
\setlist[itemize,1]{label={\textbullet},leftmargin=1.6em}
\setlist[itemize,2]{label={\textendash},leftmargin=1.4em}
\setlist[enumerate,1]{label=\arabic*., leftmargin=1.8em}

\usepackage{tabularx}
\usepackage{array}
\usepackage{booktabs}
\usepackage{longtable}

\usepackage{fancyhdr}
\usepackage{lastpage}
\pagestyle{fancy}
\fancyhf{}
\renewcommand{\headrulewidth}{0.3pt}
\renewcommand{\footrulewidth}{0pt}
\fancyhead[L]{\footnotesize\sffamily\color{lpGray}Aufgabenheft \textperiodcentered{} Kurs 22 \textperiodcentered{} Tag 12}
\fancyhead[R]{\footnotesize\sffamily\color{lpGray}Revenue Streams \& Skalierung}
\fancyfoot[L]{\footnotesize\sffamily\color{lpGray}\textcopyright{} Can Siebert}
\fancyfoot[R]{\footnotesize\sffamily\color{lpGray}\thepage{} / \pageref{LastPage}}

\fancypagestyle{plain}{%
  \fancyhf{}%
  \renewcommand{\headrulewidth}{0pt}%
  \fancyfoot[C]{\footnotesize\sffamily\color{lpGray}\thepage{} / \pageref{LastPage}}%
}

\usepackage[explicit]{titlesec}
\titleformat{\section}[block]
  {\sffamily\Large\bfseries\color{lpAccent}}
  {\thesection.}{0.6em}{#1}
  [{\vspace{2pt}\color{lpAccent}\titlerule[0.6pt]}]
\titleformat{\subsection}[block]
  {\sffamily\large\bfseries\color{lpInk}}
  {\thesubsection}{0.6em}{#1}
\titlespacing*{\section}{0pt}{14pt}{8pt}
\titlespacing*{\subsection}{0pt}{10pt}{4pt}

\usepackage[hidelinks,unicode]{hyperref}

\newcommand{\akzent}[1]{{\caveatfont\color{lpAmber}#1}}
\newcommand{\fachbegriff}[1]{\textbf{#1}}

\newcommand{\niveaubadge}[2]{%
  \tcbox[on line, colback=#1, colframe=#1, coltext=white,
         boxrule=0pt, arc=2pt, left=5pt,right=5pt,top=1pt,bottom=1pt,
         nobeforeafter]{\sffamily\bfseries\footnotesize #2}%
}
\newcommand{\nivLi}{\niveaubadge{lpL1}{L1\,\textbullet\,Wissen}}
\newcommand{\nivLii}{\niveaubadge{lpL2}{L2\,\textbullet\,Anwendung}}
\newcommand{\nivLiii}{\niveaubadge{lpL3}{L3\,\textbullet\,Management}}

\newcommand{\zeit}[1]{%
  \tcbox[on line, colback=lpGrayLight, colframe=lpGrayLight, coltext=lpInk,
         boxrule=0pt, arc=2pt, left=5pt,right=5pt,top=1pt,bottom=1pt,
         nobeforeafter]{\sffamily\footnotesize\textbf{$\sim$\,#1}}%
}

\newcommand{\aufgabenkopf}[4]{%
  \par\vspace{6pt}
  \noindent{\sffamily\Large\bfseries\color{lpAccent}Aufgabe #1 \textendash{} #2}\par
  \vspace{2pt}
  \noindent{\color{lpAccent}\rule{\linewidth}{0.6pt}}\par
  \vspace{4pt}
  \noindent #3 \hfill \zeit{#4}\par
  \vspace{6pt}
}

\newcommand{\ytquery}[1]{%
  \par\vspace{4pt}
  \noindent{\footnotesize\sffamily\color{lpGray}\textbf{Lernvideo zum Thema:}\ %
    \href{https://studyflix.de/search?query=#1}{\textcolor{lpAccent}{Studyflix-Treffer}}\ ·\ %
    \href{https://www.youtube.com/results?search\_query=#1}{\textcolor{lpAccent}{YouTube-Treffer}}\\%
    \textit{\textcolor{lpGray}{Stichworte: #1}}}\par
}

\newcommand{\antwortlinien}[1]{%
  \par\vspace{6pt}
  \foreach \x in {1,...,#1}{%
    \noindent\makebox[\linewidth]{\dotfill}\\[6pt]
  }
}
\usepackage{pgffor}

% =========================================================================
\begin{document}

\begin{titlepage}
  \pagestyle{empty}
  \vspace*{1.5cm}
  \noindent{\sffamily\footnotesize\color{lpGray}LERNPLATT \textperiodcentered{} BY CAN SIEBERT}\\[2pt]
  \noindent{\color{lpAccent}\rule{\linewidth}{1.2pt}}\\[4pt]
  \noindent{\sffamily\footnotesize\color{lpGray}AUFGABENHEFT \textperiodcentered{} MODUL 6 \textperiodcentered{} TAG 12 \textperiodcentered{} KURS 22 (Bo 143) \textperiodcentered{} 15.\,Juni 2026}

  \vspace{3cm}
  {\sffamily\fontsize{30}{34}\selectfont\bfseries\color{lpInk}Aufgabenheft\\Tag 12\par}

  \vspace{0.7cm}
  {\sffamily\large\color{lpGray}Revenue Streams optimieren,\\Multi-Channel skalieren,\\Wachstum verantwortlich steuern.\par}

  \vspace{1.5cm}
  {\caveatfont\LARGE\color{lpAmber}11 Aufgaben \textperiodcentered{} L1 \textperiodcentered{} L2 \textperiodcentered{} L3}\par

  \vfill

  \noindent\begin{minipage}{0.55\linewidth}
    {\sffamily\footnotesize\color{lpGray}Niveau-Mix:}\\
    {\sffamily\small\color{lpInk}3\,\textsf{L1} \textperiodcentered{} 5\,\textsf{L2} \textperiodcentered{} 3\,\textsf{L3}}\\[10pt]
    {\sffamily\footnotesize\color{lpGray}Bearbeitung:}\\
    {\sffamily\small\color{lpInk}Selbstbearbeitung im Lernpfad}
  \end{minipage}\hfill
  \begin{minipage}{0.4\linewidth}
    \raggedleft
    {\sffamily\footnotesize\color{lpGray}Dozent:}\\
    {\sffamily\small\color{lpInk}Can Siebert}\\[10pt]
    {\sffamily\footnotesize\color{lpGray}Begleitend:}\\
    {\sffamily\small\color{lpInk}Lehrskript \& Lösungsheft}
  \end{minipage}

  \vspace{0.5cm}
  \noindent{\color{lpAccent}\rule{\linewidth}{0.6pt}}
\end{titlepage}

\section*{So arbeiten Sie mit diesem Heft}
\thispagestyle{plain}

\noindent Dieses Aufgabenheft enthält elf Aufgaben zu Tag 12 \textendash{} \emph{Revenue Streams optimieren, Multi-Channel skalieren und Wachstum verantwortlich steuern}. Die Aufgaben sind in drei Niveaus gegliedert:

\begin{itemize}
  \item \nivLi{} \textendash{} \fachbegriff{Wissen.} Kurze Reproduktion und Einordnung. 5\,--\,10 Minuten pro Aufgabe.
  \item \nivLii{} \textendash{} \fachbegriff{Anwendung.} Transfer auf konkrete Plattformsituationen. 15\,--\,30 Minuten.
  \item \nivLiii{} \textendash{} \fachbegriff{Management.} Entscheidungsarchitektur und Verantwortung. 30\,--\,45 Minuten.
\end{itemize}

\noindent Jede Aufgabe enthält eine Niveau-Markierung, einen Zeit-Richtwert, den Aufgabentext, einen Antwort-Freiraum sowie eine \emph{YouTube-Suchquery} für vertiefendes Material.

\begin{infobox}[Hinweis]
Die Musterlösungen finden Sie im separaten \emph{Lösungsheft Tag 12}. Versuchen Sie jede Aufgabe zuerst eigenständig \textendash{} der Mehrwert entsteht in der eigenen Argumentation.
\end{infobox}

\clearpage

% =========================================================================
% L1
% =========================================================================
\section*{Niveau L1 \textendash{} Wissen \& Reproduktion}
\addcontentsline{toc}{section}{L1 -- Wissen}

% --- AUFGABE 1 -----------------------------------------------------------
% LERNPLATT_HILFSVIDEOS
\section*{Hilfsvideos zu diesem Tag}
\addcontentsline{toc}{section}{Hilfsvideos zu diesem Tag}
\thispagestyle{plain}

\noindent Die folgenden YouTube-Erkl\"arvideos vermitteln die Inhalte, die Sie zur Bearbeitung der Aufgaben ben\"otigen. Klicken Sie auf den Titel, um das Video direkt zu \"offnen; die vollst\"andige URL steht in Klammern.

\begin{description}[style=nextline,leftmargin=18pt,labelindent=0pt,itemsep=8pt]
  \item[1.~Revenue Streams im Business Model Canvas] \href{https://youtu.be/Sz9AIi8RMHQ}{\textcolor{lpAccent}{https://youtu.be/Sz9AIi8RMHQ}}\\[2pt]
  {\footnotesize\color{lpGray}(URL: \nolinkurl{https://youtu.be/Sz9AIi8RMHQ})}
  \item[2.~Skalierung von Wachstumsunternehmen — strategische Hebel] \href{https://youtu.be/eXjpapyKxYA}{\textcolor{lpAccent}{https://youtu.be/eXjpapyKxYA}}\\[2pt]
  {\footnotesize\color{lpGray}(URL: \nolinkurl{https://youtu.be/eXjpapyKxYA})}
  \item[3.~ROI im Marketing als Skalierungssteuerung] \href{https://youtu.be/vimK5KeVZQA}{\textcolor{lpAccent}{https://youtu.be/vimK5KeVZQA}}\\[2pt]
  {\footnotesize\color{lpGray}(URL: \nolinkurl{https://youtu.be/vimK5KeVZQA})}
  \item[4.~Wachstum versus Skalierung — Unterschiede im Geschäftsmodell] \href{https://youtu.be/wlraVgn7nlU}{\textcolor{lpAccent}{https://youtu.be/wlraVgn7nlU}}\\[2pt]
  {\footnotesize\color{lpGray}(URL: \nolinkurl{https://youtu.be/wlraVgn7nlU})}
\end{description}
\vspace{0.6em}
\noindent{\footnotesize\color{lpGray}\textit{Tipp:} \"Offnen Sie das Video, lassen Sie es im Hintergrund laufen und arbeiten Sie parallel an der Aufgabe.}

\clearpage

\aufgabenkopf{1}{Revenue Streams einordnen}{\nivLi}{8\,min}

\noindent Definieren Sie in jeweils einem Satz die drei Klassen von Revenue Streams, die wir in plattformbasierten Geschäftsmodellen unterscheiden:

\begin{enumerate}
  \item \textbf{Primäre Revenue Streams} (z.\,B.\ Provisionen, Transaktionsgebühren, Hauptabonnement)
  \item \textbf{Sekundäre Revenue Streams} (z.\,B.\ Premium-Funktionen, Werbung, Upgrades, Add-ons)
  \item \textbf{Strategische Revenue Streams} (z.\,B.\ Datenservices, Partnergebühren, Affiliate, Servicepakete)
\end{enumerate}

\noindent Ergänzen Sie einen vierten Satz: Warum ist es aus Steuerungssicht ein Unterschied, ob 95\,\% des Umsatzes über \emph{eine} Quelle oder über mehrere Streams zusammenkommen?

\ytquery{Revenue Streams Plattform primär sekundär strategisch Diversifikation Definition}

\antwortlinien{10}

\clearpage

% --- AUFGABE 2 -----------------------------------------------------------
\aufgabenkopf{2}{Multi-Channel-Kanäle benennen}{\nivLi}{8\,min}

\noindent Sechs typische Kanäle eines Multi-Channel-Ansatzes:

\begin{enumerate}
  \item Eigener Webshop
  \item Marktplatz / Plattform
  \item Partnervertrieb (Reseller, Wiederverkäufer)
  \item Direktvertrieb (Außendienst, Innendienst)
  \item Affiliate-Kanal
  \item E-Mail-Strecken \& digitale Kampagnen
\end{enumerate}

\noindent Beschreiben Sie für jeden Kanal in einem Satz:

\begin{itemize}
  \item Welche Kundengruppe spricht er typischerweise an?
  \item Welcher \emph{Hauptvorteil} und welches \emph{Hauptrisiko} kennzeichnen ihn (z.\,B.\ Reichweite vs.\ Abhängigkeit, Marge vs.\ Aufwand, Kontrolle vs.\ Geschwindigkeit)?
\end{itemize}

\ytquery{Multi-Channel Vertrieb Plattform Direktvertrieb Partner Affiliate Webshop}

\antwortlinien{14}

\clearpage

% --- AUFGABE 3 -----------------------------------------------------------
\aufgabenkopf{3}{KPIs für Plattform-Wachstum}{\nivLi}{8\,min}

\noindent Erläutern Sie kompakt sechs zentrale KPIs zur Steuerung plattformbasierter Wachstumsstrategien:

\begin{enumerate}
  \item \textbf{Conversion Rate}
  \item \textbf{Customer Acquisition Cost (CAC)}
  \item \textbf{Customer Lifetime Value (CLV)}
  \item \textbf{Churn-Rate} bzw.\ \textbf{Net Revenue Retention}
  \item \textbf{Wiederkaufrate}
  \item \textbf{Deckungsbeitrag je Kanal}
\end{enumerate}

\noindent Ergänzen Sie einen abschließenden Satz: Welche KPI zeigt \emph{echtes} Wachstum und welche nur \emph{Aktivität}? Warum?

\ytquery{Plattform KPIs Conversion CAC CLV Churn Net Revenue Retention Wachstum}

\antwortlinien{12}

\clearpage

% =========================================================================
% L2
% =========================================================================
\section*{Niveau L2 \textendash{} Anwendung \& Transfer}
\addcontentsline{toc}{section}{L2 -- Anwendung}

% --- AUFGABE 4 -----------------------------------------------------------
\aufgabenkopf{4}{LearnMarket Pro \textendash{} Revenue Streams analysieren}{\nivLii}{25\,min}

\begin{infobox}[Fallbeschreibung]
\textbf{Unternehmen:} ``LearnMarket Pro'' betreibt eine Plattform für berufliche Weiterbildung. Unternehmen können dort digitale Kurse, Live-Trainings und Coaching-Angebote buchen. Trainer, Akademien und Beratungsunternehmen stellen ihre Angebote auf der Plattform ein.

\smallskip
\textbf{Gegeben:}
\begin{itemize}
  \item Kundenseite: Unternehmen, Personalabteilungen, Führungskräfte.
  \item Anbieterseite: Trainer, Akademien, Beratungen.
  \item Aktuelle Einnahmequellen: 12\,\% Provision pro Buchung + kostenpflichtige Premium-Profile für Anbieter.
  \item Viele Kunden buchen nur einmalig, einige Anbieter wünschen mehr Sichtbarkeit.
  \item Große Unternehmen fragen nach Rahmenverträgen und Reporting.
  \item Die Geschäftsführung möchte die Einnahmequellen verbreitern.
\end{itemize}
\end{infobox}

\noindent\textbf{Arbeitsauftrag:}

\begin{enumerate}
  \item Benennen Sie mindestens \textbf{sechs mögliche Revenue Streams} für ``LearnMarket Pro''.
  \item Ordnen Sie jede Einnahmequelle einer Marktseite zu: \emph{Kunde}, \emph{Anbieter} oder \emph{Partner}.
  \item Beschreiben Sie für jede Einnahmequelle den möglichen \emph{Nutzen} und das \emph{Risiko} in jeweils einem Halbsatz.
  \item Entscheiden Sie, welche Einnahmequelle kurzfristig (0\,--\,6 Monate) am einfachsten umsetzbar ist.
  \item Entscheiden Sie, welche Einnahmequelle langfristig (24+ Monate) den größten \emph{strategischen} Wert haben könnte \textendash{} und begründen Sie den Unterschied.
\end{enumerate}

\ytquery{Plattform Revenue Streams Weiterbildung B2B Provision Subscription Reporting}

\antwortlinien{22}

\clearpage

% --- AUFGABE 5 -----------------------------------------------------------
\aufgabenkopf{5}{LearnMarket Pro \textendash{} Multi-Channel und KPIs}{\nivLii}{30\,min}

\begin{infobox}[Zusätzliche Informationen]
\textbf{Plattformkennzahlen:}
\begin{itemize}
  \item Monatlicher Plattformumsatz: 420.000\,\euro
  \item 70\,\% des Umsatzes über Einzelbuchungen
  \item 20\,\% über Premium-Anbieterprofile
  \item 10\,\% über Partnerempfehlungen
  \item Conversion Rate von Besucher zu Buchung: 1{,}4\,\%
  \item Wiederkaufrate nach 6 Monaten: 18\,\%
\end{itemize}

\smallskip
\textbf{Mögliche Kanäle und Ansätze:}
\begin{itemize}
  \item \textbf{A:} Ausbau des eigenen Plattformvertriebs (SEO, Content, On-Site-Conversion).
  \item \textbf{B:} Aufbau von HR-Partner-Netzwerken (Empfehlungsmodell, Provision).
  \item \textbf{C:} Direktansprache von Großkunden über LinkedIn / Account-based Marketing.
  \item \textbf{D:} E-Mail-Nurturing für Bestandskunden (Cross-Buchungen, Themenpfade).
  \item \textbf{E:} Affiliate-Kooperation mit Berufsverbänden und Karrierenetzwerken.
\end{itemize}
\end{infobox}

\noindent\textbf{Arbeitsauftrag:}

\begin{enumerate}
  \item Bewerten Sie jeden Kanal A\,--\,E auf einer Skala 1\,--\,5 in \emph{vier} Dimensionen: \textbf{Wirkung auf Conversion}, \textbf{Wirkung auf Wiederkauf}, \textbf{Marge}, \textbf{Komplexität}.
  \item Priorisieren Sie maximal \textbf{drei} Kanäle für die nächsten 6 Monate.
  \item Definieren Sie \textbf{vier KPIs}, an denen Sie den Erfolg dieser Priorisierung messen würden.
  \item Begründen Sie in 2\,--\,3 Sätzen: Wo wäre \emph{Conversion-Optimierung} wirtschaftlich wirkungsvoller als zusätzlicher Traffic?
\end{enumerate}

\medskip
\noindent\begin{tabularx}{\linewidth}{@{}l|c|c|c|c@{}}
\toprule
\textbf{Kanal} & \textbf{Conversion} & \textbf{Wiederkauf} & \textbf{Marge} & \textbf{Komplexität} \\
\midrule
A: Plattformvertrieb  & \rule{1.2em}{0.4pt} & \rule{1.2em}{0.4pt} & \rule{1.2em}{0.4pt} & \rule{1.2em}{0.4pt} \\
\midrule
B: HR-Partner         & \rule{1.2em}{0.4pt} & \rule{1.2em}{0.4pt} & \rule{1.2em}{0.4pt} & \rule{1.2em}{0.4pt} \\
\midrule
C: LinkedIn / ABM     & \rule{1.2em}{0.4pt} & \rule{1.2em}{0.4pt} & \rule{1.2em}{0.4pt} & \rule{1.2em}{0.4pt} \\
\midrule
D: E-Mail-Nurturing   & \rule{1.2em}{0.4pt} & \rule{1.2em}{0.4pt} & \rule{1.2em}{0.4pt} & \rule{1.2em}{0.4pt} \\
\midrule
E: Affiliate-Verbände & \rule{1.2em}{0.4pt} & \rule{1.2em}{0.4pt} & \rule{1.2em}{0.4pt} & \rule{1.2em}{0.4pt} \\
\bottomrule
\end{tabularx}

\ytquery{Multi-Channel KPIs Conversion Wiederkauf B2B Plattform Priorisierung}

\antwortlinien{14}

\clearpage

% --- AUFGABE 6 -----------------------------------------------------------
\aufgabenkopf{6}{Conversion-Funnel zerlegen}{\nivLii}{20\,min}

\begin{infobox}[Funnel-Daten ``LearnMarket Pro'']
Für ein durchschnittliches Quartal liegen folgende Daten vor:
\begin{itemize}
  \item Website-Besucher: 240.000
  \item Seitenaufrufe Kursdetail: 96.000 (40\,\%)
  \item Registrierte Interessenten: 14.400 (15\,\% von Kursdetail)
  \item Abgesendete Anfragen / Buchungswarenkörbe: 4.800 (33\,\%)
  \item Tatsächliche Buchungen: 3.360 (70\,\%)
  \item Wiederkauf nach 6 Monaten: 605 (18\,\%)
\end{itemize}
\end{infobox}

\noindent\textbf{Arbeitsauftrag:}

\begin{enumerate}
  \item Identifizieren Sie den \textbf{größten Conversion-Engpass} im Funnel und begründen Sie Ihre Wahl quantitativ (auf Basis der \emph{absoluten} Verluste, nicht nur der Prozentwerte).
  \item Schlagen Sie \textbf{drei konkrete Maßnahmen} vor, mit denen sich dieser Engpass adressieren ließe.
  \item Schätzen Sie für eine Ihrer Maßnahmen: Wenn diese den Engpass um 20\,\% verbessert \textendash{} wie viele zusätzliche Buchungen pro Quartal kämen ungefähr zustande?
  \item Welche \textbf{KPI} würden Sie installieren, um den Effekt sauber zu messen \textendash{} und welche \emph{Vanity-Metrik} würden Sie bewusst \emph{nicht} verwenden?
\end{enumerate}

\ytquery{Conversion Funnel Optimierung Engpass Plattform B2B Wachstum messen}

\antwortlinien{18}

\clearpage

% --- AUFGABE 7 -----------------------------------------------------------
\aufgabenkopf{7}{Kanalkonflikte und Kannibalisierung}{\nivLii}{22\,min}

\begin{infobox}[Kontext]
``LearnMarket Pro'' überlegt, parallel zum Plattformvertrieb eine eigene \emph{Direct-Sales-Einheit} für Großkunden aufzubauen \textendash{} mit individuellen Rahmenverträgen und niedrigeren Effektivpreisen ab 20.000\,\euro{} Jahresvolumen. Gleichzeitig laufen auf der Plattform öffentlich sichtbare Listenpreise.
\end{infobox}

\noindent\textbf{Arbeitsauftrag:}

\begin{enumerate}
  \item Identifizieren Sie \textbf{drei konkrete Kanalkonflikte}, die zwischen Plattform-Listenpreis und Direct-Sales-Rahmenvertrag entstehen können.
  \item Beschreiben Sie für jeden Konflikt einen \textbf{Stakeholder}, der diesen Konflikt zuerst spürt (Anbieter, Bestandskunde, Vertriebsteam, Partner, Affiliate).
  \item Entwickeln Sie \textbf{drei Steuerungsmechanismen}, mit denen sich Kanalkonflikte begrenzen lassen (z.\,B.\ Volumenstaffeln, Kanal-Caps, Preislogik, Provisionsmodell, Kanal-Verantwortlichkeiten).
  \item Formulieren Sie eine \textbf{Faustregel} in einem Satz: Wann ist Kannibalisierung akzeptabel \textendash{} und wann nicht?
\end{enumerate}

\ytquery{Kanalkonflikt Kannibalisierung Plattform Direct Sales Multi-Channel B2B}

\antwortlinien{16}

\clearpage

% --- AUFGABE 8 -----------------------------------------------------------
\aufgabenkopf{8}{Skalierung vs.\ Stabilisierung \textendash{} Streams einordnen}{\nivLii}{20\,min}

\noindent\textbf{Arbeitsauftrag:}

\noindent Eine wachsende B2B-Plattform betreibt sechs Revenue Streams. Ordnen Sie jeden Stream in genau eine der vier Kategorien ein: \fachbegriff{skalieren}, \fachbegriff{stabilisieren}, \fachbegriff{testen} oder \fachbegriff{begrenzen}. Begründen Sie Ihre Wahl jeweils in 1\,--\,2 Sätzen.

\medskip
\noindent\textbf{Streams:}
\begin{enumerate}
  \item \textbf{Transaktionsprovision} (8\,\% je Buchung) \textendash{} 70\,\% des Umsatzes, stabil, sinkende Marge wegen Preiswettbewerb.
  \item \textbf{Premium-Anbieterprofile} \textendash{} 15\,\% des Umsatzes, Anbieter beklagen unklaren Nutzen.
  \item \textbf{Unternehmens-Abo mit Reporting} (neu) \textendash{} 50.000\,\euro{} Pilotumsatz, sehr hohe Marge, aber nur 3 Pilotkunden.
  \item \textbf{Daten-/Benchmark-Service} (Idee) \textendash{} hoher strategischer Wert, Datenschutzrisiken noch ungeklärt.
  \item \textbf{Affiliate-Netzwerk} \textendash{} 8\,\% des Umsatzes, wächst organisch, Partner werden zahlreicher und schwer zu steuern.
  \item \textbf{Display-Werbung von Drittanbietern} \textendash{} 2\,\% des Umsatzes, niedrige Marge, gefährdet Nutzererlebnis.
\end{enumerate}

\noindent Schließen Sie mit \emph{einem} Vergleichssatz: Welcher Stream verdient die meiste Managementaufmerksamkeit in den nächsten 6 Monaten \textendash{} und warum?

\ytquery{Revenue Streams skalieren stabilisieren begrenzen Plattform Portfolio Wachstum}

\antwortlinien{16}

\clearpage

% =========================================================================
% L3
% =========================================================================
\section*{Niveau L3 \textendash{} Management \& Entscheidungsarchitektur}
\addcontentsline{toc}{section}{L3 -- Management}

% --- AUFGABE 9 -----------------------------------------------------------
\aufgabenkopf{9}{Fallstudie MarketSupply 360 \textendash{} 12-Monats-Architektur}{\nivLiii}{45\,min}

\begin{infobox}[Fallstudie: MarketSupply 360]
``MarketSupply 360'' ist eine B2B-Plattform, auf der mittelständische Unternehmen industrielle Verbrauchsmaterialien, Ersatzteile und Serviceleistungen beschaffen. Anbieter zahlen aktuell eine Provision pro Transaktion. Zusätzlich gibt es erste Premium-Profile und vereinzelte Partnerumsätze durch Logistikdienstleister. Die Plattform wächst, aber die Profitabilität bleibt hinter den Erwartungen zurück.

\smallskip
\textbf{Rahmenbedingungen:}
\begin{itemize}
  \item Jahresumsatz über Plattform: 18\,Mio.\,\euro
  \item Plattformprovision: durchschnittlich 8\,\%
  \item Aktive Anbieter: 1.400 \textperiodcentered{} Aktive Kunden: 9.500
  \item Conversion Rate: 1{,}9\,\% \textperiodcentered{} Wiederkauf nach 12 Mon.: 32\,\%
  \item Durchschnittlicher Warenkorb: 620\,\euro
  \item Sortimentsmix: 65\,\% Standardprodukte, 25\,\% Ersatzteile, 10\,\% Services
  \item Budget für 12 Monate: 600.000\,\euro
\end{itemize}

\smallskip
\textbf{Aktuelle Probleme:}
\begin{itemize}
  \item Starke Abhängigkeit von Transaktionsprovisionen.
  \item Premium-Profile liefern wenig messbaren Mehrwert.
  \item Kunden brechen häufig vor Anfrage oder Bestellung ab.
  \item Anbieter beklagen unklare Sichtbarkeitslogik.
  \item Partnerkanäle sind nicht systematisch aufgebaut.
  \item Management erhält viele Kennzahlen, aber keine klare Entscheidungslogik.
\end{itemize}

\smallskip
\textbf{Mögliche Ansätze:}
\begin{itemize}
  \item \textbf{A:} Unternehmensabo für Großkunden mit Reporting, Rahmenkonditionen, Priority-Support.
  \item \textbf{B:} Conversion-Optimierungs-Programm (Funnel, Suche, Produktdaten, Checkout).
  \item \textbf{C:} Partnernetzwerk mit Logistik- und Servicepartnern, neue Provisionslogik.
  \item \textbf{D:} Daten- und Benchmark-Service für Einkäufer (Marktpreise, Verfügbarkeit, Lieferzeiten).
  \item \textbf{E:} Affiliate-Modell mit Branchenmedien und Einkäuferportalen.
  \item \textbf{F:} Display-Werbung / Sponsored Listings ausbauen.
\end{itemize}
\end{infobox}

\noindent\textbf{Arbeitsauftrag:}

\begin{enumerate}
  \item Analysieren Sie die Profitabilitäts- und Wachstumssituation in 3\,--\,5 Punkten.
  \item Bewerten Sie die Ansätze A\,--\,F nach: Umsatzpotenzial \textperiodcentered{} Marge \textperiodcentered{} CLV-Wirkung \textperiodcentered{} Risiko (Kannibalisierung, Anbieterakzeptanz, Datenschutz) \textperiodcentered{} Komplexität \textperiodcentered{} Strategischer Beitrag.
  \item Entwickeln Sie eine \textbf{priorisierte 12-Monats-Architektur}: Welche Ansätze werden \emph{skaliert}, \emph{pilotiert}, \emph{stabilisiert}, \emph{zurückgestellt} oder \emph{beendet}?
  \item Verteilen Sie das Budget von 600.000\,\euro{} grob auf Ihre priorisierten Maßnahmen.
  \item Definieren Sie eine \textbf{KPI-Dashboard-Logik} mit maximal 6 KPIs, die das Management monatlich betrachten sollte.
  \item Treffen Sie \emph{mindestens eine Entscheidung unter Unsicherheit} und begründen Sie diese ausdrücklich aus Senior-Level-Perspektive.
\end{enumerate}

\ytquery{B2B Plattform Revenue Streams Skalierung CRO Mittelstand Profitabilität}

\antwortlinien{34}

\clearpage

% --- AUFGABE 10 ----------------------------------------------------------
\aufgabenkopf{10}{Kanal-Caps und Diversifikationsgrenzen}{\nivLiii}{30\,min}

\noindent\textbf{Diskussionsaufgabe.}

\noindent Multi-Channel-Diversifikation senkt Abhängigkeit \textendash{} erhöht aber Komplexität, Kannibalisierungsrisiko und Steuerungsaufwand. Senior-Management muss aktiv \emph{Grenzen} setzen: welcher Kanal darf wachsen, welcher nicht, wann ist \emph{Stoppen} besser als \emph{Skalieren}?

\medskip
\noindent\textbf{Arbeitsauftrag:}

\begin{enumerate}
  \item Beschreiben Sie in eigenen Worten (3\,--\,4 Sätze) drei \emph{Diversifikationsfallen}, in die wachsende Plattformen häufig laufen. Beispiele: Kanalproliferation ohne Marge, KPI-Inflation ohne Entscheidung, Partnerwildwuchs ohne Steuerung, Pricing-Inkonsistenz, Datenbrüche.
  \item Wählen Sie ein \textbf{konkretes Plattformbeispiel} aus Ihrer Erfahrung oder Recherche (z.\,B.\ eine SaaS-Plattform, ein B2B-Marktplatz, ein Affiliate-getriebenes Reisemodell). Erläutern Sie:
  \begin{itemize}
    \item Welche Kanal-Caps wären sinnvoll (Umsatzanteil, Margenuntergrenze, Komplexitätsbudget)?
    \item Welcher Stream sollte trotz Wachstumspotenzial bewusst \emph{begrenzt} werden \textendash{} und warum?
    \item Welche Governance-Struktur (CRO, Plattformmanagement, Controlling, Partner Management) entscheidet darüber?
  \end{itemize}
  \item Formulieren Sie eine \textbf{Faustregel} (1\,--\,2 Sätze): Ab wann muss ein neuer Revenue Stream \emph{nicht} eingeführt werden, obwohl er Umsatz brächte?
\end{enumerate}

\ytquery{Plattform Diversifikation Kanal Cap Revenue Portfolio Senior Management Grenzen}

\antwortlinien{20}

\clearpage

% --- AUFGABE 11 ----------------------------------------------------------
\aufgabenkopf{11}{Transferprojekt \textendash{} Revenue- und Performance-Architektur 24 Monate}{\nivLiii}{45\,min}

\begin{infobox}[Rolle und Ausgangslage]
\textbf{Ihre Rolle:} Chief Revenue Officer (oder Chief Platform Officer, Leiter:in Monetarisierung, Growth Manager:in, Vertriebsarchitekt:in, externe:r Strategieberater:in) eines wachsenden Plattformunternehmens.

\smallskip
\textbf{Ausgangslage:} Mehrere Revenue Streams, verschiedene Vertriebskanäle und erste Partneransätze sind vorhanden. Umsätze steigen, aber die Komplexität wächst schneller: Einige Kanäle wachsen schnell bei niedriger Marge, andere sind profitabel, aber schwer skalierbar. Partner bringen Reichweite und Abhängigkeit. Die Conversion Rate ist ausbaufähig, die KPI-Landschaft ist unübersichtlich.

\smallskip
\textbf{Rahmenbedingungen:}
\begin{itemize}
  \item Budget begrenzt.
  \item Marktinformationen unvollständig.
  \item Wettbewerber wachsen aggressiv.
  \item Einige Partner fordern exklusive Konditionen.
  \item Vertrieb, Marketing, Produkt und Controlling bewerten Erfolg unterschiedlich.
  \item Zu viele parallele Initiativen können die Organisation überfordern.
  \item Geschäftsführung erwartet schnelle Ergebnisse \emph{und} nachhaltige Skalierung.
\end{itemize}
\end{infobox}

\noindent\textbf{Arbeitsauftrag:} Entwickeln Sie eine strukturierte \emph{Entscheidungsarchitektur} \textendash{} keine reine Maßnahmenliste \textendash{} für Revenue-Optimierung, Multi-Channel-Diversifikation, Performance-Analyse und Skalierung über die nächsten 24 Monate. Erarbeiten Sie schriftlich:

\begin{enumerate}
  \item Ein \textbf{Zielbild} für die Revenue-Architektur (3\,--\,5 Sätze): Umsatz, Margenziel, Diversifikationsgrad, Steuerbarkeit.
  \item Übersicht bestehender und möglicher Revenue Streams.
  \item Bewertung der Streams nach Umsatz, Marge, CLV, Risiko, Skalierbarkeit, Kundennutzen, Komplexität.
  \item \textbf{Multi-Channel-Architektur} mit eigener Plattform, Partnerkanälen, Direktvertrieb, Affiliate und Kundenbindung.
  \item \textbf{Entscheidung}, welche Ansätze \emph{skaliert}, \emph{getestet}, \emph{stabilisiert}, \emph{zurückgestellt} oder \emph{beendet} werden.
  \item \textbf{Conversion-Optimierungslogik} entlang des gesamten Funnels (Sichtbarkeit \textrightarrow{} Klick \textrightarrow{} Anfrage \textrightarrow{} Kauf \textrightarrow{} Wiederkauf \textrightarrow{} Upgrade).
  \item \textbf{KPI-Architektur} für Managemententscheidungen: Umsatz, Deckungsbeitrag, CLV, CAC, Churn, Conversion, Wiederkauf, Partnerbeitrag, Net Revenue Retention.
  \item \textbf{Governance-Struktur} zwischen Vertrieb, Marketing, Plattformmanagement, Partner Management, Controlling, Geschäftsführung.
  \item Risikoanalyse: Kannibalisierung, Kanalabhängigkeit, Margenverlust, Partnerkonflikte, Datenbrüche, organisatorische Überforderung.
  \item Drei \textbf{Managemententscheidungen unter Unsicherheit}, die trotz unvollständiger Informationen verantwortet werden.
\end{enumerate}

\noindent\textbf{Zusatzanforderung:} Treffen Sie mindestens \emph{eine bewusste Entscheidung gegen} eine attraktive Wachstumsoption \textendash{} weil sie kurzfristig Umsatz brächte, langfristig aber Marge, Steuerbarkeit, Kundennutzen oder strategische Unabhängigkeit gefährden würde. Begründen Sie diese Entscheidung aus Senior-Level-Perspektive.

\ytquery{Chief Revenue Officer Architektur Plattform 24 Monate Senior Management Decision}

\antwortlinien{38}

\clearpage

% --- Abschluss -----------------------------------------------------------
\section*{Abschluss}
\thispagestyle{plain}

\noindent Sie haben alle elf Aufgaben bearbeitet. Vergleichen Sie nun Ihre Antworten mit dem \emph{Lösungsheft Tag 12}.

\medskip
\begin{infobox}[Was Sie jetzt können sollten]
\begin{itemize}
  \item Primäre, sekundäre und strategische Revenue Streams unterscheiden.
  \item Multi-Channel-Kanäle gegen ihre Risiken (Kannibalisierung, Komplexität, Abhängigkeit) abwägen.
  \item KPI-Ketten von Sichtbarkeit über Conversion bis Net Revenue Retention aufbauen.
  \item Conversion-Engpässe quantitativ identifizieren und gezielt adressieren.
  \item Eine Revenue-Architektur entwerfen, in der Streams \emph{skaliert}, \emph{stabilisiert}, \emph{getestet} oder \emph{begrenzt} werden.
  \item Kanal-Caps und Diversifikationsgrenzen aus Senior-Perspektive setzen.
  \item Wachstum gegen Marge, Steuerbarkeit und Kundennutzen verantwortlich abwägen.
\end{itemize}
\end{infobox}

\vspace{1cm}
\noindent{\color{lpAccent}\rule{\linewidth}{0.6pt}}\\[4pt]
{\footnotesize\sffamily\color{lpGray}Lernplatt \textperiodcentered{} by Can Siebert \textperiodcentered{} Kurs 22 (Bo 143) \textperiodcentered{} Tag 12 \textperiodcentered{} Aufgabenheft \textperiodcentered{} \textcopyright{} 2026}

\end{document}
